\subsection{Detección de Ciclo Negativo (Bellman-Ford + Ciclos en el Grafo de Padres)}

\graybold{Preguntas guía:}

\begin{itemize}
\item ¿El grafo tiene pesos negativos y me piden detectar (e imprimir) un ciclo de peso total negativo?
\item ¿Me sirve un origen virtual conectado a todos los nodos con peso 0, para encontrar ciclos negativos en cualquier parte del grafo y no solo los alcanzables desde un nodo?
\item ¿Los límites son del orden de $n \cdot m \sim 10^7$, de modo que en Python necesito cortar temprano en vez de hacer siempre las $n$ rondas?
\item ¿Debo reconstruir el ciclo en orden, no solo decir que existe?
\end{itemize}

\graybold{Utilidad:} Bellman-Ford calcula caminos mínimos con pesos negativos relajando todas las aristas en rondas, y detecta ciclos negativos: si después de $n-1$ rondas algo sigue mejorando, hay uno. Es la herramienta para grafos con pesos negativos (donde Dijkstra falla), para detectar arbitraje (ciclos de ganancia en cambio de monedas, usando $-\log$ de las tasas), y para sistemas de restricciones de diferencias $x_b - x_a \le c$, que son factibles exactamente cuando el grafo asociado no tiene ciclos negativos. La variante de esta entrada agrega una revisión periódica del grafo de padres, que en Python hace la diferencia entre terminar en milisegundos o recorrer todas las rondas.

\graybold{Complejidad:} \bigO{n m} en el peor caso (hasta $n$ rondas de $m$ aristas), más \bigO{n} por cada revisión del grafo de padres, que se hace cada 32 rondas.

\graybold{Fundamento teórico:} Inicializar todas las distancias en 0 equivale a agregar un origen virtual con aristas de peso 0 hacia cada nodo: así cualquier ciclo negativo es alcanzable y se detecta. Tras la ronda $r$, la distancia de cada nodo es a lo sumo el peso del mejor camino de $r$ aristas desde el origen virtual. Sin ciclos negativos todo camino mínimo es simple y tiene a lo sumo $n$ aristas (contando la virtual), así que tras $n-1$ rondas reales nada puede mejorar; si en la ronda $n$ algo mejora, hay un ciclo negativo. Esperar a la ronda $n$ es lo que en Python resulta caro, y el atajo es el grafo de padres: cada nodo apunta al nodo desde el que obtuvo su distancia actual. Un lema clásico (CLRS, capítulo 24) garantiza que CUALQUIER ciclo en ese grafo tiene peso negativo: al fijar el último padre del ciclo, su distancia bajó estrictamente, y sumando las desigualdades $\text{dist}[v] \ge \text{dist}[\text{padre}(v)] + w$ alrededor del ciclo se obtiene peso total $< 0$. Además, si hay un ciclo negativo, el grafo de padres termina conteniendo un ciclo. Como cada nodo tiene a lo sumo un padre, el grafo de padres es un grafo funcional y se revisa en \bigO{n}: desde cada nodo no visitado se sigue la cadena de padres marcando con el número del recorrido actual; si se llega a un nodo marcado en ESTE mismo recorrido, se encontró un ciclo, y ese nodo pertenece a él. Revisando cada 32 rondas, un ciclo negativo se detecta apenas aparece en los punteros, en vez de esperar a la ronda $n$. Si no hay ciclos negativos, se corta cuando una ronda no mejora nada. Para imprimir el ciclo, se parte de un nodo que está en él y se siguen los padres hasta volver; eso lo recorre al revés, así que se invierte.

\graybold{Ejercicio aplicado}

Dado un grafo dirigido de $n$ nodos y $m$ aristas con pesos que pueden ser negativos, decidir si contiene un ciclo de peso negativo y, si lo tiene, imprimir sus nodos en orden.

\graybold{I/O:} $n \le 2500$, $m \le 5000$ con tres enteros por arista $\Rightarrow$ lectura total + tokenización (patrón 2), armando la lista de aristas con \texttt{zip} de tres slices de paso 3. Como se acepta cualquier ciclo negativo, la verificación comprueba que la respuesta sea YES exactamente cuando existe un ciclo negativo (según Floyd-Warshall) y que el ciclo impreso cierre, use aristas existentes y sume negativo: se hizo en 800 grafos aleatorios (600 de hasta 9 nodos y 200 de 30 a 60 nodos, donde la revisión periódica sí se activa), con lazos negativos, aristas repetidas y varios ciclos, sin fallos. Con $n = 2500$ y $m = 5000$ se midieron siete grafos adversos: los que tienen ciclo negativo (de 3 nodos escondido, o de 2500 nodos) terminan en \aprox{}0.02s gracias a la revisión de padres, y el peor caso es una cadena negativa con las aristas listadas al revés, sin ciclo, que obliga a las $n-1$ rondas: \aprox{}0.53s frente a un límite de 1 segundo. Se comparó con SPFA (Bellman-Ford con cola), que es más rápido en la mayoría de los casos pero llega a \aprox{}1.00s en una cadena numerada al revés, conocida por forzar a SPFA a mejorar cada nodo de a una unidad por vuelta de la cola; por eso se eligió Bellman-Ford por rondas.

\graybold{Código solución}

\begin{lstlisting}[language=Python]
import sys

def ciclo_en_padres(n, par):
    # grafo funcional (cada nodo tiene a lo sumo un padre):
    # devuelve un nodo que esta en un ciclo, o 0 si no hay ciclos
    estado = [0]*(n + 1)              # 0 = no visto, s = en el recorrido s, -1 = terminado
    for s in range(1, n + 1):
        if estado[s]:
            continue
        x = s
        while x and not estado[x]:
            estado[x] = s
            x = par[x]
        if x and estado[x] == s:      # volvio a un nodo de ESTE recorrido: ciclo
            return x
        x = s
        while x and estado[x] == s:   # el recorrido no cerro ciclo: se marca terminado
            estado[x] = -1
            x = par[x]
    return 0

def solve():
    data = sys.stdin.buffer.read().split()
    n = int(data[0]); m = int(data[1])
    E = list(zip(map(int, data[2:2+3*m:3]), map(int, data[3:3+3*m:3]), map(int, data[4:4+3*m:3])))
    dist = [0]*(n + 1)                # todo en 0: origen virtual conectado a todos
    par = [0]*(n + 1)
    x = 0
    for ronda in range(1, n + 1):
        cambio = 0
        for a, b, c in E:
            nd = dist[a] + c
            if nd < dist[b]:
                dist[b] = nd; par[b] = a; cambio = b
        if not cambio:                # nada mejoro: no hay ciclo negativo
            break
        # todo ciclo en el grafo de padres es negativo: se busca cada 32 rondas
        if ronda == n or ronda % 32 == 0:
            x = ciclo_en_padres(n, par)
            if x or ronda == n:
                if not x:             # respaldo: mejora en la ronda n => hay ciclo;
                    x = cambio        # n pasos hacia atras caen dentro de el
                    for _ in range(n): x = par[x]
                break
    if not x:
        sys.stdout.write("NO\n"); return
    ciclo = [x]; y = par[x]
    while y != x:                     # los padres recorren el ciclo al reves
        ciclo.append(y); y = par[y]
    ciclo.append(x); ciclo.reverse()
    sys.stdout.write("YES\n" + " ".join(map(str, ciclo)) + "\n")

solve()
\end{lstlisting}
